
La centrifugación se define como la separación de sustancias de diferente densidad mediante movimiento giratorio. La base de esta operación unitaria es pues la transferencia de momento.

\section{Fundamentos de centrifugación}

\subsection{Fuerzas}

Hay dos fuerzas involucradas en la centrifugación:
- Fuerza centrífuga: No es una fuerza como tal, sino que es el resultado de la tendencia de un cuerpo a seguir en linea recta.
- Fuerza centrípeta: Es aquella fuerza que "tira" hacia el centro.

\subsection{Campos de aceleración}

Hay dos campos de fuerza que actuan sobre las partículas sometidas a una centrifugación y que aplican aceleración sobre ellas:

\begin{itemize}
    \item Gravitacional
    \item Centrifugo
\end{itemize}

\subsubsection{Gravitacional}

En primer lugar está el gravitacional, el cual es resultado de la gravedad generada por la Tierra y tira de los objetos hacia el suelo.

En este la aceleración ($a$) es la típica de $9.81\frac{m}{s^{2}}$, mientras que la fuerza impulsora es la diferencia de densidades ($\Delta\rho$) entre la partícula involucrada y el fluido.

\subsubsection{Centrifugo}

Este se origina del movimiento giratorio. Su aceleración ($a$) es:

\begin{equation}
    a = \omega^{2} \cdot r
\end{equation}

En esta ecuación $\omega$ es la velocidad de rotación en radianes, y se calcula a partir de las revoluciones ($N$) de la siguiente manera:

\begin{equation}
    \omega = \frac{2\pi \cdot N}{60}
\end{equation}

Por otro lado, $r$ es la distancia radial desde el eje de rotación a la partícula

\subsection{Ley de Stokes}

    \resumebox{definicion}{Ley de Stokes}{Cuando se aplica una fuerza a una partícula en un medio continuo ésta se acelera, hasta que alcanza una velocidad a la cual la resistencia a su movimiento iguala a la fuerza aplicada.}

Derivado de la ley de Stokes y sus ecuaciones podemos obtener que la velocidad de sedimentación de una partícula en flujo laminar ($v_{\infty}$) es:

\begin{equation}
    v_{\infty} = \frac{d_{p}^{2} \cdot \Delta\rho \cdot a}{18\mu}
\end{equation}

Donde:

\begin{itemize}
    \item $d_{p}$ es el diametro de la partícula
    \item $\Delta\rho$ es la diferencia de densidad entre la partícula y el medio en que se encuentra
    \item $a$ es la **aceleración** a la que se ve sometida
    \item $\mu$ es la viscosidad del medio
\end{itemize}

Esta aceleración puede ser alguna de las dos vistas anteriormente (gravitacional o centrífuga), tras lo cual al reemplazar tenemos una ecuación para cada caso:

\begin{align}
    \text{Campo gravitacional} &\quad\longrightarrow\quad v_{g}   &= \frac{d_{p}^{2} \cdot \Delta\rho \cdot (g)}{18\mu}
    \\\\
    \text{Campo centrífugo} &\quad\longrightarrow\quad v_{\omega} &= \frac{d_{p}^{2} \cdot \Delta\rho \cdot (\omega^{2}\cdot r)}{18\mu}
\end{align}








\section{Diseño de centrífugas}

\subsection{Tubular}

La ecuación utilizada para el diseño de centrífugas tubulares es:

\begin{equation}
    Q = (v_{g}) \cdot \left[ \left(\frac{\pi\cdot\omega^{2}\cdot L}{g}\right) \cdot \left(\frac{R_{0}^{2} - R_{1}^{2}}{\ln\left(\frac{R_{0}}{R_{1}}\right)}\right) \right]
\end{equation}

Donde:

\begin{itemize}
    \item $v_{g}$ es la velocidad de precipitación por acción del campo gravitacional
    \item $\omega$ es la velocidad de rotación en radianes
    \item $L$ es el la longitud del eje central de la centrífuga
    \item $g$ es la aceleración de la gravedad
    \item $R_{0}$ es la distancia del centro de la centrífuga a sus paredes internas
    \item $R_{1}$ es la distancia del centro de la centrífuga a las paredes de la interfase líquida
\end{itemize}

De esta ecuación podemos identificar varios términos, los cuales se ven mejor si se colorea la ecuación:

\begin{equation}
    Q = \textcolor{red}{(v_{g})} \cdot \left[ \textcolor{blue}{\left(\frac{\pi\cdot\omega^{2}\cdot L}{g}\right)} \cdot \textcolor{green}{\left(\frac{R_{0}^{2} - R_{1}^{2}}{\ln\left(\frac{R_{0}}{R_{1}}\right)}\right)} \right]
\end{equation}

Aquí los términos azul y verde que están entre corchetes se designan en su conjunto como sigma mayúscula ($\Sigma$):

\begin{equation}
    \Sigma = \left[ \left(\frac{\pi\cdot\omega^{2}\cdot L}{g}\right) \cdot \left(\frac{R_{0}^{2} - R_{1}^{2}}{\ln\left(\frac{R_{0}}{R_{1}}\right)}\right) \right]
\end{equation}

En el caso del término azul este es adimensional, esto porque:

\begin{itemize}
    \item $\pi$ no tiene dimensiones
    \item Las  dimensiones de $\omega$ son $\frac{1}{s}$
    \item Las  dimensiones de $L$ son $m$
    \item Las  dimensiones de $g$ son $\frac{m}{s^{2}}$
\end{itemize}

Y reemplazando:

\begin{equation}
    \frac{\pi\cdot\omega^{2}\cdot L}{g} \quad\longrightarrow\quad \frac{\left(\frac{1}{s}\right)^{2} \cdot m}{\frac{m}{s^{2}}} \quad\longrightarrow\quad \frac{\frac{m}{s^{2}}}{\frac{m}{s^{2}}} \quad\longrightarrow\quad \text{Se cancela todo}
\end{equation}

Por otro lado, el segundo término, el verde, unicamente tiene el valor de $R$, cuyas dimensiones son $m$. Reemplazando tenemos:

\begin{equation}
    \frac{R_{0}^{2} - R_{1}^{2}}{\ln(\frac{R_{0}}{R_{1}})} \quad\longrightarrow\quad \frac{m^{2} - m^{2}}{\ln(\frac{m}{m})} \quad\longrightarrow\quad \text{ln cancela unidades} \quad\longrightarrow\quad m^{2}
\end{equation}

Con todo esto, tenemos basícamente que $\Sigma$ tiene unidades de $m^{2}$, por lo que representa esencialente el área de contacto entre fluido y centrífuga.

Todo este analisis dimensional lo que nos revela es que los dos términos principales de la ecuación, $v_{g}$ y $\Sigma$, representan cada uno una parte fundamental del proceso de centrifugación:

\begin{itemize}
    \item $v_{g}$ describe la velocidad con que las partículas sedimentan, y revela características de tales partículas y del fluido, tales como su densidad o viscosidad.
    \item $\Sigma$ describe características de la centrifuga, tales como el área de contacto.
\end{itemize}

Podemos observar por último el término en azul también:

\begin{equation}
    \left(\frac{\pi\cdot\omega^{2}\cdot L}{g}\right)
\end{equation}

Este nos describe cuanto va a aumentar la velocidad de sedimentación ($v_{g}$) tras aplicar la centrifugación. Esto porque una partícula no sedimenta a la misma velocidad estando unicamente sujeta a un campo gravitacional, que cuando además está sujeta también a un campo centrífugo.

\subsubsection{$Q_{50\%}$}

Si bien la ecuación anterior nos sirve para diseñar un proceso de centrifugación para una partícula de diametro dado, estos cálculos quedan "justos" para tal proceso.  Una mejor alternativa puede ser diseñar el proceso para que centrifuge el 50% de las partículas de una diametro más pequeño, lo cual a su vez aseguraría centrifugar el 100% de las partículas más grandes, o sea, las "partículas objetivo".

\begin{equation}
    Q_{50\%} = (2v_{g}) \cdot \left[ \left(\frac{\pi\cdot\omega^{2}\cdot L}{g}\right) \cdot \left(\frac{R_{0}^{2} - R_{1}^{2}}{\ln\left(\frac{2R_{0}^{2}}{R_{0}^{2} + R_{1}^{2}}\right)}\right) \right]
\end{equation}



\subsection{De discos}

Para una centrífuga de discos la ecuación para calcular el flujo es la siguiente:

\begin{equation}
    Q = (v_{g}) \cdot \left[ \left(\frac{2\pi\cdot\omega^{2}\cdot n}{3g}\right) \cdot \left(\frac{R_{0}^{3} - R_{1}^{3}}{\tan(\theta)} \right) \right]
\end{equation}

Donde, además de las variables ya mencionadas en la ecuación de centrifugas tubulares, tenemos que:

\begin{itemize}
    \item $n$ son los números de discos
    \item $\theta$ es el ángulo de inclinación de tales discos
    \item $R_{0}$ ahora es el radio interno de los discos
    \item $R_{1}$ ahora es el radio externo de los discos
\end{itemize}